\section{材料与方法}
\subsection{研究设计、观测数据与环境覆盖}
本研究使用 Ortiz-Andrade 和 Rojas-Perez\cite{ortiz2026}公开的派生数据与脚本，以及本项目已有实验输出。研究覆盖波多黎各，保留采用定点或行进协议、调查时长为 5--240~min、行进距离不超过 5~km 的完整 eBird 清单。响应值为 1 表示报告该物种，为 0 表示未报告；未报告不等于已确认缺失。

生境（habitat）指物种生活的环境。本研究以海拔和植被绿度描述部分生境条件，以清单记录和模型估计的占域状态分析鸟类与这些条件的关系。文中的生境利用为这一关系的概括，不表示直接测量了觅食、繁殖行为或生境质量。记录层面的再次记录，指同一地点在后续时期又有清单报告该物种；它同时受到物种存在和探测机会影响，与模型估计的占域持续概率分别报告。

分析时段为 2013 年 1 月 1 日至 2025 年 12 月 31 日。时期划分沿用原文定义。María 前截至 2017 年 9 月 19 日；两次飓风之间为 2017 年 9 月 20 日至 2022 年 9 月 17 日；Fiona 后从 2022 年 9 月 18 日开始。这些分类描述观测相对于风暴的时间位置，没有量化局地风暴暴露，也未提供未受扰动的反事实。因此，时期对比是时间关联，而非飓风因果效应的估计。

海拔来自 SRTM\cite{farr2007}，NDVI 表示植被绿度\cite{tucker1979}。公开的 Landsat 流程在 $\pm16$ 天内选择最近的有效观测，并将影像限定于对应时期。本研究使用已发布的协变量，不重新进行卫星提取。全部分析包含 \nChecklists{} 份清单，其中 \nComplete{} 份具有 NDVI，\nMissingNDVI{} 份（\nMissingPct{}\%）缺失。环境 GLM 和主要常规模型基准使用完整子集。我们按时期及报告结果描述缺失情况，不预设缺失机制为随机。这一组成检查参考了公众科学趋势估计中空间与时间采样偏差相互作用的研究证据\cite{backstrom2025}。原有 EVI2 比较作为植被指数的次级检查\cite{jiang2008}。

图~\ref{fig:area}将清单覆盖与验证区块置于同一景观中。图中 24,876 个地点来自环境提取前的坐标表；环境导出表因坐标精度差异包含 24,856 个坐标对。两者均不同于占域地点数及验证区块数。采样位置代表实际观测机会，并非空间均匀布设的调查。
\begin{figure}[htbp]\centering
\incfig[0.94\linewidth]{fig01_study_area_and_spatial_blocks}
\caption{波多黎各地形、清单覆盖（上图）及一次空间折分配（下图）。完整的 $0.08^\circ$ 区块被分配至不同折；相邻训练区块与测试区块之间没有设置距离缓冲。}\label{fig:area}
\end{figure}

\subsection{数据核验与条件报告模型}
\label{sec:verify}
文件哈希、样本计数、筛选条件和拟合系数用于检查数据与模型复现。两条日期--坐标重复记录均予以保留，以保持与原研究的样本一致。
\ifzh
表~\ref{tab:screen}列出研究对象的来源筛选。参考研究的实施适合度得分为 5.1，高于其余候选的 3.4--4.1，主要优势是同时提供观测数据、环境协变量、分析脚本及校验信息。这使本研究能够从同一输入出发，连续检查报告格局、占域过程和预测验证，减少不同资料来源造成的比较差异。
\else
Table~\ref{tab:screen} summarizes source selection. The reference study scored 5.1 for implementation suitability, compared with 3.4--4.1 for the other candidates. Its combination of observations, environmental covariates, analysis scripts and checksums enabled a linked evaluation of reporting patterns, occupancy processes and predictive validation from a common input.
\fi
\inctab{screen}
\ifzh
数据核验通过 53/54 项检查，GLM 的 44/44 项检查均在规定容差内通过，11 个占域系数及标准误也得到复现（图~\ref{fig:verify}；表~\ref{tab:verify}）。唯一的数据检查差异涉及两条日期--坐标组合重复产生的多余记录，约占全部清单的 0.0022\%，且均未报告目标物种。该结果将数值差异的排查重点从文件版本和样本筛选转向模型内部的观测协变量对应关系。
\else
Data verification passed 53 of 54 checks; all 44 GLM checks met their tolerances, and all 11 occupancy coefficients and standard errors were reproduced (Figure~\ref{fig:verify}; Table~\ref{tab:verify}). The sole data-check discrepancy involved two excess records with duplicated date--coordinate combinations, approximately 0.0022\% of all checklists, both without a species report. This localized the subsequent investigation to observation-covariate alignment within the model rather than differences in input versions or filtering.
\fi
\begin{figure}[htbp]\centering\incfig{fig02_verification_and_reproduction}
\caption{\ifzh
数据核验与模型数值复现。颜色区分检查类别，系数比较使用原模型排列。
\else
Data verification and numerical model reproduction. Colors distinguish check categories; coefficient comparisons use the original model ordering.
\fi}\label{fig:verify}\end{figure}
\FloatBarrier
{\footnotesize\renewcommand{\arraystretch}{0.96}\inctab{verify}}
\FloatBarrier


五个二项 GLM 从仅含截距依次加入时期、观察努力、环境及时期--海拔交互，全部使用相同的 NDVI 完整子集。原文选定模型为
\begin{equation}
\operatorname{logit}(q_i)=\alpha+\beta_{P(i)}+\beta_E E_i+\delta_{P(i)}E_i+\beta_V V_i+\beta_D D_i+\beta_L L_i,
\label{eq:glm}
\end{equation}
其中 $q_i$ 为报告概率，$E_i$ 为以 100~m 为单位的海拔，$V_i$ 为 NDVI，$D_i$ 为时长（分钟），$L_i$ 为距离（千米）；María 前为参照时期。已有 AIC、模型权重与伪 $R^2$ 用于概括候选模型的拟合支持度\cite{burnham2002,mcfadden1974}，与留出预测表现分别报告。

绘制海拔曲线时，将 NDVI、时长及距离固定在完整样本均值。100 和 900~m 的预测用于概括低、高海拔对比，不作为生物学阈值。沿用公开分析，在 0--1300~m、间隔 10~m 的均匀网格上计算
\begin{equation}
C_t=\frac{\sum_h h\widehat q_t(h)}{\sum_h\widehat q_t(h)}.
\end{equation}
这一概率加权海拔质心描述拟合曲线，不是种群数量、地理分布中心或个体移动的估计。年度报告率区间沿用已有二项区间，未处理观察者或地理依赖。

\subsection{占域持续性与观测模型校正}
\label{sec:ordering}
动态模型以坐标四舍五入形成的 $0.01^\circ$ 网格作为地点，将三个时期作为主要调查时期。纳入地点须在每个时期具有清单。在每个地点--时期内，先按年份、再按原始行序排列，保留前三次记录；不足三次的位置记为缺失。最终 \nOccSites{} 行、九列的检测历史包含 \nOccasions{} 次已观测调查及 \nOccMissing{} 个缺失位置。地点协变量在筛选访问次数之前，使用全部符合条件记录计算，并在保留地点间标准化。因此，NDVI 是跨时期地点汇总值，不是同一地点植被恢复的时间轨迹。

选定的 \texttt{unmarked} 模型\cite{fiske2011}以海拔与 NDVI 预测初始占域（$\psi$）和局地灭绝（$\varepsilon$），以 NDVI 预测定殖（$\gamma$），以时期预测探测（$p$）。时期未单独进入转移分量，因此该模型不估计某次飓风特有的灭绝或定殖变化。

令 $z_{g,t}$ 表示网格 $g$ 在主要时期 $t$ 的潜在占域状态，占域持续概率（occupancy persistence probability）定义为 $\Pr(z_{g,t+1}=1\mid z_{g,t}=1)=1-\varepsilon_g$，即前一时期已占据的网格在下一时期仍被占据的概率\cite{mackenzie2003}。它不表示同一个体持续停留、个体存活率或整个时期内不间断的出现。定殖概率 $\gamma$ 以此前未占据为条件，而观测到的首次记录或再次记录不直接判定定殖。文中的占域持续性均指上述网格尺度的状态转移。

观测协变量应按地点优先顺序与检测矩阵对应，即一个地点的九次观测排列在下一地点之前。公开脚本却把每个调查位置的时期标签重复到全部地点。已有带标签的小矩阵检查及受控重拟合仅改变这一排列，保持响应、地点协变量、公式及优化设置不变。我们比较系数与不确定性，以 $\widehat\beta\pm1.96\,\mathrm{SE}$ 计算近似 95\% Wald 区间。这是选定模型内部的敏感性检验，并未重新选择全部校正候选模型。主要时期长达数年，也对时期内状态封闭假设提出较强要求；其转移估计不是年度种群动态速率。

\subsection{留出地理区块中的报告预测}
\label{sec:protocols}
预测目标仍为清单层报告。十一个常规配置包括 GLM m3、GLM m4、样条 GLM、加入地形的 m4、随机森林、LightGBM、XGBoost、多层感知机（MLP），以及后三者加入地形的版本。树模型提供灵活的表格数据预测\cite{breiman2001,ke2017,chen2016}。其输入包括环境、努力量、时期、季节项及坐标；原文 GLM 保持已发表的协变量集合。地形配置另加入 20 项 SRTM 描述变量，概括海拔、坡度、坡向和四个尺度的邻域变化。因此，比较评价的是实际模型与信息的组合，并未单独隔离架构效应。

保存的配置中，随机森林使用 400 棵树、最小叶节点样本数为五；两类提升树均使用 600 个估计器及 0.05 的学习率，LightGBM 为 63 个叶节点，XGBoost 最大深度为六；MLP 包含三个宽度为 128 的隐藏层。MLP 标准化和样条节点等需估计的变换在训练折内学习。这些设置描述已有实验，本次不依据报告的测试结果重新调参。

主要验证把边长 $0.08^\circ$（约 8--9~km）的完整地理区块分配至十个样本量均衡折，使用三个区块原点重复，十一个配置共享全部 30 个测试折。每折可以由不相邻区块组成，训练与测试区块交界处没有距离缓冲。该方案检验指定的岛内地理留出情境，不代表任意距离外推。次级随机验证采用三次分层五折划分。随机与空间方案的训练比例和测试组成随各自划分规则确定。原有时间实验以三个种子重复预测同一批 Fiona 后清单，用于评价较早时期训练关系对末期报告的预测。地理与时间留出分别对应不同的推广任务，符合按预测用途安排验证的原则\cite{spacetime2025}。

主要指标为 AUC、平均精确率（AP）、Brier 得分，以及排序前 10\% 清单对正报告的召回率\cite{saito2015,brier1950}，分别反映总体判别、正报告排序、概率误差和固定排序预算下的覆盖。模型比较使用不加权折均值及折间标准差，并报告绝对差值、相对改善和前十分位正报告覆盖。重复划分中的训练集相互重叠，因而不把各折视为独立重复进行显著性分类。合并预测曲线另采用十折预测导出。ROC 和精确率--召回率评价排序，12 个等频组绘制可靠性曲线，15 个等频组计算期望校准误差。尺度敏感性分别使用 $0.04^\circ$、$0.08^\circ$、$0.16^\circ$ 和 $0.32^\circ$ 区块，EVI2 替换实验采用六折。

\subsection{\ifzh 环境分层、过程概率与监测预算实验\else Environmental stratification, process probabilities and monitoring budgets\fi}
\label{sec:envmethods}
\ifzh
补充分析将海拔分为低于 300~m、300--不足 600~m 和不低于 600~m 三层；NDVI 采用 68,976 条完整清单的三分位切点 0.4172 和 0.7006，形成九个环境组合。这些分层用于探索环境异质性，不作为物种的生境阈值。占域分析将相同 NDVI 切点应用于地点跨时期均值，并报告各组合的地点数；少于十个地点的组合单独标记。
\else
Supplementary analyses stratified elevation as below 300~m, 300--below 600~m, and at least 600~m. NDVI tertile cutpoints of 0.4172 and 0.7006, computed from 68,976 complete checklists, produced nine environmental combinations. These exploratory strata describe heterogeneity rather than species habitat thresholds. Occupancy summaries apply the same NDVI cutpoints to cross-period site means, report site counts, and flag combinations with fewer than ten sites.
\fi
\ifzh
NDVI 缺失比较先汇总各海拔层的报告率与缺失率，再拟合缺失指示变量与报告之间的二项关联模型，依次加入时期和海拔样条，以及调查时长的对数、距离的 $\log(1+x)$ 变换、季度和经纬度二次项。标准误按 214 个 $0.08^\circ$ 空间区块聚类。模型使用信赖域优化获得稳定起点，并检查收敛、梯度及协方差。报告回升使用复现的 m4 公式重新拟合，将每个环境组合内全部记录的协变量分布固定，分别替换时期后求平均预测。Fiona 后相对 María 前的差值区间通过聚类协方差和 delta 法计算；原始分层报告率另进行 1,000 次空间区块自助抽样。
\else
NDVI-missing comparisons first summarize report and missingness rates within elevation bands. Binomial reporting models then relate reports to a missingness indicator, sequentially adding period and an elevation spline, followed by log duration, $\log(1+x)$ distance, quarter, and quadratic geographic-coordinate terms. Standard errors cluster 214 geographic blocks of width $0.08^\circ$. Trust-region optimization supplies stable starting values, with convergence, gradient and covariance checks. Reporting rebound is standardized using a refit of the reproduced m4 formula. Each environmental cell retains its pooled covariate distribution while period is replaced in turn. Post-Fiona minus Pre-Mar\'ia intervals use clustered covariance and the delta method; raw stratified rates additionally use 1,000 spatial block bootstrap samples.
\fi
\ifzh
校正占域模型在 DGX 上重新拟合，系数与已有校正结果的最大差异小于 $10^{-4}$。对每个地点预测初始占域、占域持续概率 $1-\varepsilon$ 和定殖 $\gamma$，再按环境组合平均。2,000 次联合正态参数模拟保留完整协方差，给出概率与对比的 95\% 区间。固定海拔 100~m、NDVI 从 0.30 增至 0.76 的对比用于分离低海拔内的绿度响应；这些取值近似对应低海拔低、高绿度地点的中位数。另从九列检测历史中筛选相邻两个时期各有三次访问的地点--时期对，统计前期已报告后再次报告的比例。此比例描述等访问次数下的重复报告，与潜在占域转移分别呈现。
\else
The corrected occupancy model was refitted on the DGX, reproducing existing corrected coefficients within $10^{-4}$. Site predictions of initial occupancy, occupancy persistence probability $1-\varepsilon$, and colonization $\gamma$ were averaged by environmental cell. Two thousand joint normal coefficient draws preserve the full covariance and yield 95\% probability and contrast intervals. A contrast at 100~m from NDVI 0.30 to 0.76 isolates the fitted greenness response; these values approximate the low- and high-greenness medians among low-elevation sites. Adjacent period pairs with three visits in each period were also extracted from the nine-column detection history to summarize repeated reports conditional on a previous report. This equal-visit reporting proportion is presented separately from latent occupancy transitions.
\fi
\ifzh
新增随机森林采用原 400 棵树、最小叶节点样本数五和种子零，在一次十折空间划分中重新训练。逐折核验事件标识符和地理区块的训练--测试交集均为空，并重新计算 m4 的首折预测验证原预测文件的行对应关系。分层预测比较包括 m4、随机森林和两种地形提升树。监测预算在每个测试折固定为清单数的 10\% 向下取整。全局方案直接选该折最高分记录，分层方案按各海拔层样本量分配约 10\%，采用最大余数分配保持每折总名额完全相同。两方案均选择 6,890 条清单。报告 AP、Brier、各层入选比例及正报告覆盖率；覆盖率差的区间采用 1,000 次空间区块自助抽样，条件于本次划分的拟合预测。该实验评价已有清单的排序与环境覆盖，不把清单数直接换算为实地访问成本。
\else
A new random forest used the original 400 trees, minimum leaf size five and seed zero in one ten-fold spatial partition. Every fold was checked for disjoint training/test event identifiers and geographic blocks; recomputing the first m4 fold validated the row alignment of the saved predictions. Stratified comparisons include m4, random forest and both terrain-augmented boosting models. Each test fold receives a budget equal to the floor of 10\% of its checklist count. Global ranking selects the highest scores within that fold; stratified ranking allocates approximately 10\% within each elevation band, using largest remainders to preserve exactly the same fold budget. Both select 6,890 checklists. AP, Brier, selection fractions and positive-report coverage are reported by stratum. Coverage-contrast intervals use 1,000 spatial block bootstrap samples conditional on this partition's fitted predictions. The experiment evaluates record ranking and environmental coverage; checklist counts are not converted into field-visit costs.
\fi

\subsection{神经预测结构与训练来源诊断}
\label{sec:audit}
原 HERON 流程将报告得分分为两个有界分支得分 $q_i=s_i d_i$，再与常规模型组合。我们重建已有分折规则，用调查事件标识符匹配完整记录表与全部记录表。原流程独立划分两表，却把相同折编号视为相同留出记录。本研究先检查测试记录与神经训练来源是否重叠，再解释性能；内部堆叠也存在类似的独立分配。诊断依据原训练划分和观测标识符计算重叠比例。

此外，分组损失使用 $1-\prod_i(1-s_i d_i)$，仅依赖分支乘积，因而自身不能把某个分支辨识为占域或探测。共享地点状态的占域模型对应的概率是 $\Pr(\text{至少一次报告})=\psi_g[1-\prod_j(1-p_{gj})]$。输入限制与努力量单调性不证明当前损失等价于这一共享状态似然。因此，神经环境与努力分支均以拟合得分解释，占域持续性由动态占域模型的转移分量评价。神经组件采用 20 折消融，堆叠采用六折消融，直接配置比较采用十折；各对比在自身折集合内分析。
\begin{figure}[htbp]\centering\incfig[\linewidth]{主图/1}
\caption{\ifzh
环境分析与监测建议概览。（a）eBird 记录、SRTM 海拔、NDVI 和调查协变量用于条件报告、动态占域与随机森林报告预测三类互补分析。（b）分析结果为固定山地复访、低地绿度对照及环境分层内候选记录排序提供依据，以指导后续调查。
\else
Overview of the environmental analyses and monitoring implications. (a) eBird records, SRTM elevation, NDVI and survey covariates inform complementary analyses of conditional reporting, dynamic occupancy and random forest reporting prediction. (b) The findings motivate fixed upland revisits, lowland greenness contrasts and within-stratum candidate ranking to guide follow-up surveys.
\fi}\label{fig:arch}
\end{figure}
\FloatBarrier

\FloatBarrier
